\subsection{Monomial Orders}  
What do we need to do to carry over the long division and gcd algorithms to multivariate setting? 

In the case of on variable, our monomials were naturally ordered by degree. The leading term was defined based on this order and long division doesn't need anything but 
the access to the leading term function and the operations + and $\cdot$ on $k[x]$. \\[0.3em]
Say, how to order $x^2y, \ xy^2$ and $xyz$? 

The answer is: there is no unique way to do that. In the multivariate vase, there are many ways to order monomials. Still, we need to define what a reasonable ordering is. 
Let's first clarify what if means to order a set.  

\begin{definition}[title = Definition 2.1.1]
  Consider a binary relation $\preccurlyeq$ (\textit{\textcolor{airforceblue}{curly $\leq$}}) on a set $M$. We say that $\preccurlyeq$ is a non-strict total order on $M$ if the following hold:
  \begin{itemize}
    \item Antisymmetry: \[x \succcurlyeq y, y \succcurlyeq x \Rightarrow x = y, \quad \forall x,y \in M\] 
    \item Transitivity: \[x \succcurlyeq y, y \succcurlyeq z \Rightarrow x \succcurlyeq z, \quad \forall x,y,z \in M\]
    \item Totality: \[x \succcurlyeq y \text{ or } y \succcurlyeq x, \quad \forall x,y \in M\]
  \end{itemize} 
\end{definition}  

\begin{center}
  \begin{addmargin}[0.6cm]{0.6cm}
      \textit{\textcolor{airforceblue}{For a finite $M$ it is like listing the elements of $M$ in some order in an array, but for an infinite $M$, such an orering is not 
      necessarily resembling an infinite array. Eg.}}
    \[\underbrace{\bullet \cdot \cdot \cdot \cdot \cdot \cdot \cdot}_{\mathclap{\substack{\text{infinitely} \\ \text{many}}}} \ \ \underbrace{\bullet \cdot \cdot \cdot \cdot \cdot \cdot \cdot}_{\mathclap{\substack{\text{infinitely} \\ \text{many}}}}\ \ \underbrace{\bullet \cdot \cdot \cdot \cdot \cdot \cdot \cdot}_{\mathclap{\substack{\text{infinitely} \\ \text{many}}}} \dots\]
     \textit{\textcolor{airforceblue}{The above picture sketches a situation where a linear ordering is an infinite concatenation of infinite arrays.}}
  \end{addmargin}
\end{center}

Having a non-strict total order $\preccurlyeq$ on $M$, we can also consier the respective strict total order $\prec$ given by 
\[x \prec y \Leftrightarrow x \preccurlyeq y, x \neq y\]
Clearly, $\preccurlyeq$ determines $\prec$ and, vice versa, $\prec$ determines $\preccurlyeq$. So, it is enough to descibe one of them when we totally orer some given set $M$. \\

We want to order the set of all monomials in \[x = (x_1, \ldots, x_n): \quad \{x^{\alpha} \mid \alpha \in \mathbb{Z}^n_{\geq 0}\}\] 

Monomial are determined by their exponent vector. This means, we want to totally order the set $\mathbb{Z}^n_{\geq 0}$ of all exponent vectors. \\[0.3em] 
For $n = 1$, our ordering was: 
\[x^0 \prec x^1 \prec x^2 \prec x^3 \prec \cdots\]
corresponding to $1 < 2< 3< \dots$. 

The leading term LT(f) of a univariate polynomial $f \in k[x]\backslash\{0\}$ is the term of the larges monimal with respect to the above order. The long division algorithm 
in the univariate case multiplies polynomials with terms. 

Therefore, it relies on the following property: \[x_{\alpha} \succ x^{\beta} \Rightarrow x^{\alpha} \cdot x^{\gamma} \succ x^{ \beta} \cdot x^{\gamma} \quad ( \alpha, \beta, \gamma \in \mathbb{Z}_{\geq 0}) \] 
That means, the order is preserved under muliplication with any monimial. \\[0.3em]
There is yet another crucial property: 

When we carry out the long division the degree of the remainder $r$ in progress gets smaller in each iteration and the ermination criterion $\deg \ r < \deg \ g$ is 
achieved after finitely many iterations. For the monomial orders in the multivariate case, we need a condition that would guarantee \underline{finite termination}.
The above discussion provides motivation for the following definition\dots


\begin{definition}[title = Definition 2.1.2]
    A strict monomial order $\succ$ on $\mathbb{Z}^n_{\geq 0}$ is binary relation $\mathbb{Z}^n_{\geq 0}$ such that: 
    \begin{itemize}
      \item[(i)] $\succ$ is astrict total order on $\mathbb{Z}^n_{\geq 0}$ 
      \item[(ii)] If $\alpha, \beta, \gamma \in \mathbb{Z}^n_{\geq 0}$, then $\alpha \succ \beta$ implies $\alpha + \gamma \succ \beta + \gamma$  
      \item[(iii)] Every non-empty subset  $A$ of $\mathbb{Z}^n_{\geq 0}$ has the minimal element with respect to $\succ$, which is an elemend $\beta \in A$ satifying $\alpha \succ \beta$ for all $\alpha \in A, \alpha \neq \beta$. 
    \end{itemize}
\end{definition}
 
By identifying $\mathbb{Z}^n_{\geq 0}$ with $\{x^{\alpha}\mid \alpha \in \mathbb{Z}^n_{\geq 0}\}$ we also speak about monomial orders on $\{x^{\alpha}\mid \alpha \in \mathbb{Z}^n_{\geq 0}\}$. 

\begin{remark}[title = Remarks]
  \begin{itemize}
    \item[(i)] guarantess that any two monomials are comparable 
    \item[(ii)] guarantees compatibility with the multiplication by a monomial 
    \item[(iii)] guarantees that the multivariate long division (to be introduced) terminates. ONly infinite sets $A$ matter in (iii), as finite one have the minimum with respect to any linear order.  
  \end{itemize}
\end{remark}

\begin{remark}[title = Remark]
  A set with a total order is called totally ordered (\textit{\textcolor{airforceblue}{surprise!! :)}}). It's up to use to choose if you want to define the strict or non-strict one.
\end{remark}

\begin{definition}[title = Definition 2.1.3]
  A descending chain of a totally ordered set $(M, \succ)$ is a sequence of element that satisfy $a_0 \succ a_1 \succ a_2 \succ \cdots$. A chain can be finite of countably infinite.
\end{definition}

\begin{lemma}[title = Lemma 2.1.4]
  Let $(M, \succ)$ be a totally ordered set. The the following are equivalent: 
  \begin{itemize}
    \item[(i)] Every descending chain of $M$ is finite.
    \item[(ii)] Every non-empty subset $A$ of $M$ has the minimal element with respect to the order $\succ$.  
  \end{itemize}

  \paragraph{Proof:} It's enough to verify the contraposition $\overline{(i)} \Leftrightarrow \overline{(ii)}$ \\[0.3em] 
  If some non-semtpy set $A \subseteq M$ has no minimal element with respect to $\succ$, we can pick any $a_0 \in A$. Since $a_0$ is not minimal in $A$, there exits $a_1 \in A$ with $a_0 \succ a_1$. 
  And so on to infinity. We get the infinite chain $a_0 \succ a_1 \succ a_2 \succ \cdots$ of elements. \\[0.3em]
  Conversely, if $a_0 \succ a_1 \succ a_2 \succ \cdots$ is an infinite chain, then the set $\{a_i\mid i \in \mathbb{N}_0\}$ has no minimal element because for every $a_i$, the 
  element $a_{i +1}$ is even smaller. 
    \begin{flushright}
      $\Box $
    \end{flushright}
\end{lemma}

\clearpage

\begin{center}
  \begin{addmargin}[0.6cm]{0.6cm}
    \textit{\textcolor{airforceblue}{There are infinitely many choices of different monomial orders. We probably don't have the time to consider all of them :). So,
    let's stick to the following three that turn out to be rather useful.}}
  \end{addmargin}
\end{center}

\begin{definition}[title = Definition 2.1.5]
  {
  \centering
  \begin{tabular}{r l}
      lex = & lexicographic order \\ 
      grlex = & graded lexicographic order (deglex is another common abbreviation) \\
      grevlex = & graded reverse lexicographic order
  \end{tabular}
  } \\[0.3em]

  For $\alpha = ( \alpha_1, \ldots, \alpha_n), \beta = (\beta_1, \ldots, \beta_n) \in \mathbb{Z}^n_{\geq 0}$

  \begin{align*}
      \alpha \succ_{lex} \beta \quad :\Leftrightarrow \quad & \alpha_1 = \beta_1, \ldots, \alpha_{i-1} = \beta_{i-1}, \alpha_i > \beta_i \text{ for some } i \in \{1, \ldots, n\} \\
      \intertext{\hspace*{2cm} \textit{\textcolor{airforceblue}{dictionary order: in out dictionary all words have length $n$ and the alphabet is $\mathbb{Z}_{\geq 0}$} }}
      \alpha \succ_{grlex} \beta \quad :\Leftrightarrow \quad & |\alpha| > |\beta| \text{ or } |\alpha| = |\beta|, \alpha \succ_{lex} \beta \\
      \intertext{\hspace*{2cm} \textit{\textcolor{airforceblue}{order by total degree, ot equal tie break by lex}}} 
      \alpha \succ_{grevlex} \beta \quad :\Leftrightarrow \quad & |\alpha| > |\beta| \text{ or } \textcolor{magenta}{ \alpha_i < \beta_i}, \alpha_{i+1} = \beta_{i+1}, \ldots, \alpha_n = \beta_n \text{ for some } i \in \{1, \ldots, n\} 
    \intertext{\hspace*{2cm} \textit{\textcolor{airforceblue}{order by total degree, if equal} \textcolor{magenta}{unusual} \textcolor{airforceblue}{tie break.}}} 
  \end{align*}
\end{definition}

\begin{remark}[title = Remark]
  lex and grlex look rather natural (\textit{\textcolor{airforceblue}{right?!?}}). lex can be used for elemination. glex has a better control of the number of iterations in the long division.
  grevlex is even besser in this respect. 
\end{remark}

\vspace{0.3em} 
\begin{center}
  \textcolor{gray}{------------------------------------------- 09.11.26 --------------------------------------------} \\[0.3em]
\end{center}
\vspace{0.3em} 

\[
  \langle g_1, \ldots, g_s \rangle = \langle \underbrace{gcd(g_1, \ldots, g_s)}_{\text{definition of gcd in terms of ideals}} \rangle 
\]

\begin{proposition}[title = Proposition 2.1.6]
    $\succ_{lex}, \succ_{grlex}$ and $\succ_{grevlex}$ are all monomial orders on $\mathbb{Z}^n_{\geq 0}$

    \paragraph{Proof:} 
    We will do the proof for $\succ_{lex}$ and leae the two others as exercises. \\[0.3em]
    Clearly, $\succ_{lex}$ is a total order. 

    For $\alpha, \beta, \gamma \in \mathbb{Z}^n_{\geq 0}$, one has $\alpha \succ_{lex} \beta \quad \Rightarrow \alpha + \gamma \succ_{lex} \beta + \gamma$. It remains t show 
    that $\succ_{lex}$ is a well ordering. Consider any nonßempty subset $A$ of $\mathbb{Z}^n_{\geq 0}$.  

    Let's consider \[A_1 = \{\alpha_1 \mid (\alpha_1, \ldots, \alpha_n) \in A\} \subseteq \mathbb{Z}_{\geq 0}\]
    Let $\beta_1 = \min A_1$. 

    Let's introduce \[ A_2 := \{\alpha_2 \mid (\beta_1, \alpha_2, \ldots, \alpha_n) \in A\} \subseteq \mathbb{Z}_{\geq 0}\] Let $\beta_2 = \min A_2$ and so on 
    iteratively. 
    \[A_t =\{ \alpha_t \mid (\underbrace{beta_1, \ldots, \beta_{t-1}}_{fixed}, \alpha_t, \ldots, \alpha_n) \in A\}, \beta_t = \min A_t\quad (t = 1, \ldots, n)\]
    By construction, $(\beta_1, \ldots, \beta_n) = \min A$. 

    \begin{flushright}
      $\Box $
    \end{flushright}
\end{proposition}

\begin{remark}[title = Remark]
    For each of the three orders above, there are $n!$ versions of it, depending on the sorting of the vraiables $x_1, \ldots, x_n$. 
\end{remark}

\begin{definition}[title = Definition 2.1.7]
    Let $f \in k[x_1, \dots, x_n]\backslash\{0\}, \ f = \sum_{\alpha} c_{\alpha} x^{\alpha}$ and let $\succ$ be a monomial order on $\mathbb{Z}^n_{\geq 0}$. 
    With respect to this order, we have the following notions:
    
    \begin{itemize}
      \item[(i)] $mdeg\ f$, the multidegree of $f$, is the maximum $\alpha$ over all $\alpha$ such that $c_{\alpha} = 0$. 
      \item[(ii)] $LC(f) = c_{mdeg\ f}$ is the leading coefficient. 
      \item[(iii)] $LM(f)x^{mdeg\ f}$ is the leading monomial  
      \item[(iv)] $LT(f) = LC(f) \cdot LM(f) = c_{mdeg\ f} \cdot x^{mdeg\ f}$ is the leading term.
    \end{itemize}
    We used the convention $mdeg(0) = - \infty \prec \alpha$ (for $\alpha \in \mathbb{Z}^n_{\geq 0}$).
\end{definition}

In what follows, we will assume that an underlying monomial order has been fixed. 

\begin{lemma}[title = Lemma 2.1.8]
    Let $f, g \in k[x_1, \dots, x_n]\backslash\{0\}$. Then 
    \begin{itemize}
      \item[(i)] $mdeg (f \cdot g) = mdeg(f) + mdeg(g)$
      \item[(ii)] $mdeg (f + g) \preceq \max\{mdeg(f), mdeg(g)\}$
      \item[(iii)] $mdeg(f) \neq mdeg(g) \Rightarrow mdeg(f + g) = \max\{mdeg\ f, mdeg g \}$  
    \end{itemize} 

    \paragraph{Proof:} Straight forward. 
\end{lemma}
